\begin{frame}{역사 2: 부트스트랩(1993) $\to$ 배깅 $\to$ Breiman(2001)}
  \begin{itemize}
    \item 골튼의 원리는 약 100년 뒤, \kb{부트스트랩}(bootstrap,
      Efron \& Tibshirani 1993)이라는 도구로 돌아옴:
      \textbf{복원추출}로 원본과 같은 크기의 샘플을 계속 뽑으면 그
      통계량 분포가 \textbf{진짜 분포를 근사}
    \item 랜덤 포레스트는 이 부트스트랩을 \textbf{트리의 학습 데이터}에
      그대로 적용 — 각 트리가 ``다르게 보게 만드는 것'' =
      \kb{배깅}(B\"uhlmann \& H\"usler 1997)
    \item 마지막 조각, \textbf{특징 무작위화}는 2001년 Breiman의
      ``Random Forests''에서 완성:
      \textbf{배깅만으론}(특정 강한 특징이 있으면) 트리들이 너무
      닮는다는 문제를, ``각 분기에서 \textbf{특징의 무작위 하위 집합}만
      고려하자''는 두 번째 무작위성으로 풀었다
  \end{itemize}
  \vskip0.6em
  {\footnotesize (부트스트랩 + 특징 무작위화)의 \textbf{조합}이 바로 우리가
  쓰는 랜덤 포레스트다.}
\end{frame}
